% chapters/06_conclusion.tex - 第6章 结论与建议

\chapter{结论与建议}

\section{研究结论}
\label{sec:conclusions}

本研究围绕自研移动端计算机视觉系统SportSci Pro的测量效度验证及其在青年男性下肢抗阻训练数智化监控中的应用, 采用解释性序贯混合方法设计, 通过两项研究系统回应了第1章提出的四项核心研究目的——工具效度（目的1）、干预可行性（目的2）、干预效果（目的3）与用户接受度及决策机制（目的4）。形成如下六点主要结论。

\textbf{结论一（目的1）}：移动端计算机视觉技术具备在真实训练场景中替代专业线性位置传感器的测量学基础, 但在轻负荷热身场景中存在稳定的算法口径差异, 需通过工程校正弥补。研究一在受控实验条件下证实SportSci Pro与GymAware在杠铃后蹲平均向心速度测量上达到优秀一致性（ICC=0.940, MAE=0.047 m/s, Bias=+0.002 m/s）。研究二在真实训练场景的Rep级配对（n=43, 3名受试者）进一步获得ICC=0.991（Bias=+0.0019 m/s）, 热身级配对（n=88, AI辅助组全体11名受试者）ICC=0.987, 但存在+0.0212 m/s的稳定系统性正向偏移（88对配对100\%均为正偏）, 需在系统层实施线性校正后方可应用于负荷--速度曲线的精确建模。

\textbf{结论二（目的2）}：基于速度损失阈值的数智化监控方案在依从受试者中具备扎实的实施可行性与安全性, 但招募保留率与系统可用性两项指标未达预设推进标准, 在完成上述改进并通过可行性再评估之前, 不建议直接进入大规模确证性RCT。在PP样本中, 训练完成率达98.4\%（189/192）、高依从率达100\%（24/24）、AI辅助组速度目标命中率达80.2\%$\pm$6.9\%, 自动减组忠实度达100\%, 严重不良事件为0例, 9项预设推进标准中7项达标。然而, 主要后测完成率66.7\%（24/36）与AI辅助组SUS评分65.00$\pm$4.61分两项未达标。

\textbf{结论三（目的3）}：在等外部剂量条件下, 数智化监控与规范化RIR自我指导诱发了相似的内部负荷反应与最大力量适应, 但负荷组织形态呈现有意义的行为学分化。两组在总外部负荷（$p$=0.854）、深蹲主项负荷（$p$=0.988）、全期平均sRPE（6.5 vs. 6.4）及个体Hooper均值（20.4 vs. 19.8）上高度对等。AI辅助组在总完成重复次数上显著少于自我指导组（100.3 vs. 105.5次, $p$=0.035）, 派生的每次重复平均机械负荷高出约6.1\%, 体现了速度损失阈值即时截断的行为学特征。Hooper主观恢复状态的LMM分析显示组别与交互效应均不显著（$p$分别为0.421与0.945）, 课次主效应显著（$p$=0.004）, 两组均呈现先升后降的时间演变形态, 从内部负荷层面证实了两种反馈路径在等外部剂量下诱发相似应激反应的"等应激现象"。

\textbf{结论四（目的3）}：本研究4周干预后, 两组在深蹲绝对1RM、相对1RM与蹲跳SJ高度上未表现确证性组间优效（$p$分别为0.266、0.273、0.622；$g$分别为0.46、0.39、0.43）；在反向运动跳CMJ高度上捕获到中到大效应的方向性信号（调整后差值+2.16 cm, $p$=0.051, $g$=0.85, $g$ 95\% CI：0.02至1.69）, FDR校正后$p$=0.126, 提示数智化监控对改善牵张缩短循环爆发力可能具有方向性优势, 该结论受限于探索性研究定位, 需大样本正式试验以CMJ为主要结局指标进行验证。

\textbf{结论五（目的3）}：数智化监控在训练自我效能上产生了极显著的组间优势（调整后差值+9.88分, $p<0.001$, Hedges'$g$=2.70）, 是本研究最具一致性与稳健性的量化发现, 该效应需在三重审慎的框架下解读——开放标签期望效应混杂、模型残差非正态假设及量表天花板效应风险——具体应对策略见6.2.1节。

\textbf{结论六（目的4）}：移动端系统的用户接受度呈现典型的原型系统二元特征：功能价值获得高度认同（感知有用性、信任度、使用意愿三项TAM维度均分均超过4.4分）, 但过程性操作摩擦导致可用性评分未达标（SUS 65.00分, 仅P004一人达70分）。质性访谈揭示5类行为特征模式（积极接受型、安全依赖型、批判接受型、规则内化型、负荷导向型）, 显示不同训练经验水平的用户对自动化控制的接受度存在显著差异, 从自我决定理论视角论证了面向大众部署时"分层自适应授权"而非"一刀切自动化"的必要性。

\section{研究建议}
\label{sec:recommendations}

基于上述六点结论, 本研究对后续正式随机对照试验、数字化训练系统工程迭代与大众抗阻训练实践三个层面分别提出针对性建议。

\subsection{对后续正式随机对照试验的建议}
\label{sec:rct_recommendation}

后续正式随机对照试验应以CMJ高度作为主要结局指标（本研究观测到的中到大效应$g$=0.85提供了先验参数）, 每组约20至24人（考虑20\%脱落率后每组随机化约24人, 共48人）；若加入注意力等量对照组以控制开放标签设计带来的期望效应混杂, 则总样本量需扩展至约72人。建议在期中分析时根据累积数据进行样本量重新估算。同时, 通过压缩基线导入流程至单日集中完成, 建立主动随访机制、灵活的补测安排以及阶段性里程碑激励等综合策略, 系统性提升受试者保留率至$\geq$80\%预设阈值。

\subsection{对数字化训练系统工程迭代的建议}
\label{sec:system_recommendation}

当前原型系统的迭代重心应从功能扩展转向交互体验优化, 具体优先级依次为：极简UI与数据可视化重构（以颜色区间条替代MCV、VL\%、$\Delta$v等数值）、自动化蓝牙设备连接（将就绪时间压缩至$<$30秒）、算法反馈的可解释性增强（改为自然语言解释式输出）以及热身阶段偏移的工程层校正模块（内置App校准值=App实测值$-$0.021 m/s）。在长远设计上, 建议实现"入门--进阶--专家"三级自适应交互模式, 动态调整系统自动化程度与用户自主权限。

\subsection{对大众抗阻训练实践的建议}
\label{sec:practice_recommendation}

面向大众推广时, 不建议采用"一刀切"策略。对训练经验不足1年的初学者, 本研究质性证据表明其易在追求大重量过程中累积疲劳与技术变形, 引入基于速度损失阈值的客观反馈具有明确的应用价值。对训练年限在1至3年的进阶训练者, 两种反馈路径均能支持有效适应, 工具定位应为数据参考而非强制指令, 保留用户对最终负荷决策的主动权。对训练年限超过3年的高水平训练者, 系统应定位为高精度数据仪表盘而非决策指挥官, 充分尊重用户的自主性与本体感觉判断。

\section{研究展望}
\label{sec:future_directions}

本研究为移动端计算机视觉技术在抗阻训练监控领域的应用提供了初步的可行性证据与效度参数, 但仍有若干关键科学问题有待未来研究系统回答。

从技术维度, 当前研究聚焦于杠铃后蹲单动作场景, 未来研究应探索多动作（卧推、硬拉等）、多设备（不同品牌智能手机）与多变场景（家庭、健身房、户外）的测量学泛化性, 并结合深度学习算法迭代实现更低计算资源消耗与更强环境鲁棒性。

从应用维度, 长期（8至12周以上）干预研究有待检验速度损失阈值反馈对肌肉肥大、力量持久性与运动专项表现的长周期效应, 并追踪用户是否存在反馈依赖风险。本研究纳入对象为18至30岁具有抗阻训练经验的青年男性, 与大众健身人群及女性训练者在身体成分、激素环境和训练经验上存在系统性差异。规范化自我指导与数智化监控在女性, 不同训练水平（初学者vs.高水平）及不同健康状态（代谢性疾病、关节损伤康复期等）人群中的相对效果, 尚需在专门设计的人群特异性研究中予以验证, 不能将本研究结论直接推广至上述人群。

从方法学维度, 三臂混合方法RCT设计将为因果推断提供更强的证据等级, 并在量化实验基础上充分整合质性过程性评价, 构建数智化训练工具效果评估的完整方法学体系。
